\chapter{The classical mod--2 Adams spectral sequence}
\label{chap:classical-adams}

This chapter defines the Steenrod, Ext, tower, convergence, product, and named-class interfaces.  Concrete published differentials and low-line calculations occur only in the later external-results chapter.

\section{Steenrod algebra, Ext, towers, and convergence}

% SOURCE: KIP126/Main/Axiom/Literature/MainPaper/main.tex:126-150; PROJECT_BOUNDARY.md, Steenrod algebra and Ext.
\begin{definition}[Dual Steenrod algebra]
  \label{def:dual-steenrod-algebra}
  \uses{def:mathlib-f2}
  \notready
  The project provides the graded connected Hopf algebra
  $\mathcal A_*=H_*(H\mathbb F_2;\mathbb F_2)$, its augmentation and coproduct,
  and the graded-dual relation with the mod--$2$ Steenrod algebra $\mathcal A$.
  Mathlib v4.32.2 has no Steenrod-algebra implementation that supplies this
  structure.
  The explicit Milnor coalgebra below supplies the algebraic coproduct
  slice; its construction does not complete the Hopf structure or identify
  it with the homology of $H\mathbb F_2$.
\end{definition}

% SOURCE: KIP126/Main/Axiom/Literature/MainPaper/main.tex:126-150 and every Adams E_2-page computation.
\begin{definition}[Graded Steenrod-comodule category]
  \label{def:steenrod-comodule-category}
  \uses{def:dual-steenrod-algebra,def:mathlib-abelian,def:hf2-homology}
  \notready
  Graded right $\mathcal A_*$-comodules form the abelian category used to model
  mod--$2$ Adams input.  The interface includes degree-preserving comodule
  maps, shifts, kernels, cokernels, tensor products, the trivial comodule
  $\mathbb F_2$, and the homology comodule $H_*(X;\mathbb F_2)$.
\end{definition}

% SOURCE: KIP126/Main/Axiom/Literature/MainPaper/main.tex:126-150; internal and cohomological Ext gradings.
\begin{definition}[Bigraded Adams Ext]
  \label{def:adams-bigraded-ext}
  \uses{def:steenrod-comodule-category,def:mathlib-derived-ext}
  \notready
  The Adams group $\operatorname{Ext}_{\mathcal A}^{s,t}(M,\mathbb F_2)$ has
  derived degree $s$ and internal degree $t$.  It is represented equivalently
  by internal graded Ext in $\mathcal A_*$-comodules, with an explicit proof
  translating the paper's module convention to the comodule convention.  The
  interface supplies internal-degree shifts, external tensor pairings, and the
  action of
  $\operatorname{Ext}_{\mathcal A}^{*,*}(\mathbb F_2,\mathbb F_2)$.
  It supplies an internal product on
  $\operatorname{Ext}_{\mathcal A}^{*,*}(M,\mathbb F_2)$ only when $M$ carries
  the additional coalgebra data that induces one; there is no such product for
  an arbitrary $M$.  Mathlib's singly graded derived Ext alone does not provide
  these bigraded structures.
\end{definition}

\begin{definition}[Graded vector spaces with Cauchy tensor]
  \label{def:graded-vector-space-cauchy}
  \lean{KIP126.Algebra.GradedVectorSpace.GrVect,KIP126.Algebra.GradedVectorSpace.tensorInclusion,KIP126.Algebra.GradedVectorSpace.discreteEquivalence}
  \uses{def:mathlib-f2,def:mathlib-abelian}
  \leanok
  For a field $K$, let $\operatorname{GrVect}_K$ be the category of
  integer-graded $K$-vector spaces and degree-preserving linear maps.
  Its tensor is the Cauchy tensor
  $(V\otimes W)_n=\bigoplus_{i+j=n}V_i\otimes_K W_j$, with the actual
  coproduct inclusions. Its linear operations are pointwise on morphisms;
  the abelian structure is transported along the equivalence with functors
  from the discrete category of integers. This does not use a pointwise
  tensor product.
\end{definition}

\begin{definition}[The fixed coaugmented Milnor coalgebra]
  \label{def:fixed-milnor-graded-coalgebra}
  \lean{KIP126.Steenrod.Milnor.Coalgebra.dualSteenrodGraded,KIP126.Steenrod.Milnor.Coalgebra.comultiplication,KIP126.Steenrod.Milnor.Coalgebra.counit,KIP126.Steenrod.Milnor.Coalgebra.dualSteenrod,KIP126.Steenrod.Milnor.Coalgebra.coaugmentation}
  \uses{def:graded-vector-space-cauchy}
  The degree-$n$ carrier has the full Milnor monomials of weight $n$ as
  its basis, including the constant monomial in degree zero; negative
  degrees have empty monomial bases. In polynomial coordinates its maps are
  \[
    \Delta(\xi_k)=\sum_{i=0}^{k}\xi_{k-i}^{2^i}\otimes\xi_i,
    \qquad \xi_0=1,\qquad |\xi_k|=2^k-1.
  \]
  The counit takes the constant coefficient, and the coaugmentation is
  the constant polynomial. The coproduct is constructed from the existing
  polynomial slot-splitting operation and the corresponding Cauchy
  summands. Counitality, coassociativity and coaugmentation compatibility
  remain property proof obligations. No unrelated coalgebra maps are
  selected, and no full Hopf or topological identification follows here.
\end{definition}

\begin{definition}[Actual graded right comodules and internal shifts]
  \label{def:actual-graded-right-comodules}
  \lean{KIP126.Algebra.GradedComodule.rightTensorComonad,KIP126.Algebra.GradedComodule.RightComodule,KIP126.Algebra.GradedComodule.forget}
  \lean{KIP126.Algebra.GradedComodule.linear,KIP126.Algebra.GradedComodule.abelian,KIP126.Algebra.GradedComodule.internalShift,KIP126.Algebra.GradedComodule.trivialAt}
  \uses{def:graded-vector-space-cauchy,def:fixed-milnor-graded-coalgebra}
  For a graded coalgebra $C$, right comodules are the Eilenberg--Moore
  coalgebras of the actual comonad $V\mapsto V\otimes C$.
  Addition is the existing coalgebra-morphism addition and scalar
  multiplication is that of the underlying graded linear maps.
  Finite limits and colimits use the forgetful functor's creation
  constructions; the abelian structure is assembled from these and the
  coimage-to-image isomorphism property. The Cauchy tensor's additivity,
  linearity and finite-limit preservation, and the comparison-isomorphism
  property, are stated but not yet proved.
  Left tensoring by $K[t]$ defines the internal shift, with old degree
  $n-t$ contributing to degree $n$. A specified coaugmentation defines
  the trivial comodule $K[t]$ in every integer degree. This left tensor
  uses a graded vector space; a tensor product of two arbitrary
  $C$-comodules would additionally require suitable multiplication on $C$.
  General spectrum-homology coactions are not constructed by this node.
\end{definition}

\begin{definition}[Derived Ext in the fixed comodule category]
  \label{def:milnor-derived-ext}
  \lean{KIP126.Algebra.GradedComodule.BigradedExt,KIP126.Algebra.GradedComodule.CoefficientExt}
  \lean{KIP126.Steenrod.Milnor.Ext.RightComodule,KIP126.Steenrod.Milnor.Ext.trivialAt,KIP126.Steenrod.Milnor.Ext.SphereExt,KIP126.Steenrod.Milnor.Ext.unit}
  \uses{def:actual-graded-right-comodules,def:mathlib-derived-ext}
  In the actual graded-comodule category define
  $\operatorname{Ext}^{s,t}(M,N)=\operatorname{Ext}^s(M[t],N)$ using
  Mathlib's derived category, localized at quasi-isomorphisms.
  For the fixed Milnor coalgebra the sphere convention is
  $\operatorname{SphereExt}^{s,t}=\operatorname{Ext}^s(\mathbb F_2[t],\mathbb F_2)$,
  with $s\in\mathbb N$ and $t\in\mathbb Z$; its $(0,0)$ unit is the
  identity Ext class. This is derived Ext, not a renamed cobar quotient.
  The abelian-category property debt persists, and comparisons with cobar
  cohomology and the paper's left-module convention are separate claims.
\end{definition}

\begin{definition}[Convolution augmentation and left-module Ext]
  \label{def:milnor-left-module-ext}
  \lean{KIP126.Algebra.GradedDual.convolutionAlgebra,KIP126.Algebra.GradedDual.convolutionAugmentation}
  \lean{KIP126.Algebra.GradedModule.leftTensorMonad,KIP126.Algebra.GradedModule.LeftModule,KIP126.Algebra.GradedModule.abelian,KIP126.Algebra.GradedModule.CoefficientExt}
  \lean{KIP126.Steenrod.Milnor.Module.steenrod,KIP126.Steenrod.Milnor.Module.augmentation,KIP126.Steenrod.Milnor.Module.trivialAt,KIP126.Steenrod.Milnor.Module.SphereExt}
  \uses{def:fixed-milnor-graded-coalgebra,def:graded-vector-space-cauchy,def:mathlib-derived-ext}
  Let $A_n=\operatorname{Hom}_{\mathbb F_2}(C_n,\mathbb F_2)$ for the
  fixed Milnor coalgebra $C$. The convolution product transposes the
  ordered component $C_{i+j}\to C_i\otimes C_j$; its unit transposes the
  counit. The augmentation evaluates a degree-zero functional on the
  specified coaugmentation's value at $1$, and is zero in other degrees.
  Left modules are algebras of the actual monad $V\mapsto A\otimes V$.
  Their linear and abelian structures use the underlying maps and the
  Eilenberg--Moore creation of finite limits and colimits. Restriction along
  the augmentation gives the actual trivial degree line $k[t]$. Define
  the left-module sphere groups as $\operatorname{Ext}_A^s(k,k[t])$ in
  Mathlib's derived category. The monad, algebra and abelian structure
  property proofs remain obligations; these objects are not additional
  external inputs or new stage-boundary fields.
\end{definition}

\begin{definition}[Recursive conjugation and same-degree dualization]
  \label{def:milnor-contravariant-dualization}
  \lean{KIP126.Steenrod.Milnor.Coalgebra.Antipode.generatorImage,KIP126.Steenrod.Milnor.Coalgebra.Antipode.polynomialMap,KIP126.Steenrod.Milnor.Coalgebra.Antipode.map}
  \lean{KIP126.Steenrod.Milnor.Module.conjugation,KIP126.Steenrod.Milnor.Module.dualToLeftModule,KIP126.Steenrod.Milnor.Module.trivialDualIso}
  \uses{def:milnor-left-module-ext,def:actual-graded-right-comodules}
  For the fixed coproduct convention, construct the polynomial antipode by
  \[
    S_0=1,\qquad
    S_n=\xi_n+\sum_{1\leq i<n}\xi_{n-i}^{2^i}S_i\quad(n>0).
  \]
  Polynomial substitution extends this to all monomials, and its finite
  support restricts it to the same homogeneous carriers. Transpose this
  map to obtain $\chi:A\to A$. For a right comodule $M$, evaluation on
  its coaction first defines the right action $\varphi\cdot a$ on its
  same-degree dual. Use $a\cdot\varphi=\varphi\cdot\chi(a)$ to define
  the left action. Degreewise transpose on maps gives the actual functor
  $D:(\operatorname{Comod}_C)^{\mathrm{op}}\to\operatorname{Mod}_A$.
  Evaluation on the specified generator identifies $D(k[t])$ with $k[t]$;
  the degree is not negated. Antipode and action properties are stated but
  unproved. No equivalence of unrestricted module and comodule categories
  is asserted.
\end{definition}

\begin{definition}[Prescribed normalized cofree resolution]
  \label{def:milnor-cofree-resolution}
  \lean{KIP126.Steenrod.Milnor.Ext.Cofree.wordSpace,KIP126.Steenrod.Milnor.Ext.Cofree.term,KIP126.Steenrod.Milnor.Ext.Cofree.termPolynomial,KIP126.Steenrod.Milnor.Ext.Cofree.rawDifferential}
  \lean{KIP126.Steenrod.Milnor.Ext.Cofree.termPolynomial_injective}
  \lean{KIP126.Steenrod.Milnor.Ext.CobarResolution,KIP126.Steenrod.Milnor.Ext.CobarResolution.representativePolynomial}
  \uses{def:actual-graded-right-comodules,def:milnor-derived-ext}
  Write $W_s$ for the graded vector space on normalized Milnor words of
  length $s$. The prescribed term is the actual cofree right comodule
  $W_s\otimes C$, with the last, full coalgebra factor retaining constants.
  Polynomial coordinates concatenate the normalized word with this final
  factor, degree by degree across the Cauchy direct sum.
  Injectivity of these polynomial coordinates is a separate, currently
  unproved statement.
  A cobar resolution is an injective resolution of $\mathbb F_2[0]$ with
  comodule isomorphisms from its terms to these specified cofree terms.
  In these coordinates its differential is
  \[
    d_{\mathrm{cofree}}^s=\operatorname{insertLeft}_{s+1}
       +\sum_{j=0}^{s}\operatorname{splitSlot}_j.
  \]
  There is no right-insertion term in this resolution differential.
  The augmentation takes the actual generator $1$ to the constant
  polynomial $1$. Evaluating a map out of $\mathbb F_2[t]$ on its
  degree-$t$ generator defines its representative polynomial.
  Existence, injectivity of the terms, and the resolution comparison
  remain obligations; declaring this presentation does not prove them.
\end{definition}

\begin{theorem}[Cobar comparison with actual derived comodule Ext]
  \label{thm:cobar-derived-comodule-ext-comparison}
  \lean{KIP126.Challenge2.CobarDerivedExtComparison,KIP126.Challenge2.CobarDerivedExtComparison.internalEquiv}
  \lean{KIP126.Interface.Solution.cobarDerivedExtComparison}
  \uses{def:milnor-cofree-resolution,def:milnor-derived-ext,def:milnor-cobar-cohomology}
  For the fixed $H\mathbb F_2$ input and Milnor cooperation comparison,
  there exist one prescribed cobar resolution $R$ and $\mathbb F_2$-linear
  equivalences
  \[
    e_{s,t}:H^{s,t}_{\mathrm{cobar}}\cong
       \operatorname{SphereExt}^{s,t}\qquad(s,t\in\mathbb N).
  \]
  For every normalized cocycle $x$ of bidegree $(s,t)$, require a comodule
  map $f:\mathbb F_2[t]\to R^s$ with $fd=0$ such that
  \[
    \operatorname{representativePolynomial}_R(f)
       =\operatorname{insertRight}_s(x),\qquad
    e_{s,t}([x])=R.\operatorname{extMk}(f).
  \]
  Right insertion appends the constant final coalgebra factor. These
  representative equations bind the equivalence to the specified
  resolution in every degree. Composing with the existing cobar-to-internal
  $E_2$ comparison yields the displayed internal $E_2$-to-derived-Ext
  equivalence without choosing new coordinates.
  The target Ext is defined in all integer internal degrees, whereas this
  comparison currently quantifies nonnegative internal degrees.
  The producer theorem remains unproved. The field is correlated with
  the same foundation and Milnor input in Challenge~2; neither the
  left Steenrod-module comparison nor compatibility with Yoneda products
  follows from this additive comparison.
\end{theorem}

\begin{theorem}[Coefficient Ext comparison with prescribed representatives]
  \label{thm:milnor-coefficient-module-comparison}
  \lean{KIP126.Steenrod.Milnor.Module.dualCobarResolution}
  \lean{KIP126.Steenrod.Milnor.Module.PreservesDualCobarRepresentatives,KIP126.Steenrod.Milnor.Module.exists_dualCobarComparison}
  \uses{def:milnor-cofree-resolution,def:milnor-derived-ext,def:milnor-contravariant-dualization}
  Given a prescribed cofree resolution $R$, form the actual chain complex
  $P_s=D(R^s)$ with differential $D(d_R)$ and augmentation $D(\iota_R)$
  followed by $D(k[0])\cong k[0]$. Its projectivity and resolution
  properties remain proof obligations for these specified terms and maps.
  There exists a family of linear equivalences
  \[
    e_{s,t}:\operatorname{Ext}_C^s(k[t],k)
       \cong\operatorname{Ext}_A^s(k,k[t]),
    \qquad s\in\mathbb N,\quad t\in\mathbb Z,
  \]
  such that, for every $f:k[t]\to R^s$ with $fd_R=0$,
  \[
    e_{s,t}\bigl(R.\operatorname{extMk}(f)\bigr)
      =P.\operatorname{extMk}\bigl(D(f)\,;\,(D(k[t])\cong k[t])\bigr).
  \]
  The target map must be a cocycle for this actual dual differential.
  The all-representatives condition determines the comparison; its
  existence is stated but not proved. This is a fixed algebraic coefficient
  comparison in Def, not an external result or an additional Challenge
  field. It does not supply the general-spectrum coaction bridge or
  multiplication compatibility.
\end{theorem}

% SOURCE: KIP126/Main/Axiom/Literature/MainPaper/main.tex:126-150; finite-type standing assumption at lines 253-263.
\begin{proposition}[Module--comodule Adams Ext comparison]
  \label{prop:adams-module-comodule-ext-comparison}
  \uses{def:adams-bigraded-ext,def:admissible-spectrum,def:hf2-homology,def:unital-multiplication-data}
  \notready
  For an admissible spectrum $X$, degreewise finite duality identifies the
  paper's module expression
  \[
    \operatorname{Ext}_{\mathcal A}^{s,t}
      (H^*(X;\mathbb F_2),\mathbb F_2)
  \]
  with the corresponding internal Ext, equivalently Cotor, expression built
  from the graded $\mathcal A_*$-comodule $H_*(X;\mathbb F_2)$.  The
  isomorphism preserves $(s,t)$, external pairings, the sphere-Ext action,
  suspension, and maps induced by spectra.  When $X$ has unital multiplication
  data, it also preserves the resulting internal product.
\end{proposition}

\begin{proof}
  Use finite type to dualize each graded piece, identify Steenrod-module
  actions with dual-Steenrod coactions, and compare the bar and cobar
  resolutions degreewise.  Verify the internal grading, external tensor
  pairing, and sphere action on the chain-level comparison before passing to
  cohomology.  For a spectrum with multiplication, compare the internal
  product obtained by composing with its multiplication map.
\end{proof}

% SOURCE: KIP126/Main/Axiom/Literature/MainPaper/main.tex:134-150; PROJECT_BOUNDARY.md, Steenrod algebra and Ext.
\begin{definition}[Mod--2 Steenrod algebra and Ext]
  \label{def:steenrod-ext}
  \uses{def:adams-bigraded-ext,prop:adams-module-comodule-ext-comparison,def:hf2-homology,def:unital-multiplication-data}
  \notready
  Let $\mathcal A$ be the mod--$2$ Steenrod algebra and let
  $\mathbb F_2$ be its trivial module.  The project interface provides the
  bigraded groups
  $\operatorname{Ext}_{\mathcal A}^{s,t}(H^*(X;\mathbb F_2),\mathbb F_2)$,
  external pairings, their action by the sphere group
  $\operatorname{Ext}_{\mathcal A}^{*,*}(\mathbb F_2,\mathbb F_2)$,
  suspension/regrading maps, and functoriality in $X$.  An internal product is
  present only after supplying multiplication data on $X$; in particular it
  is canonical for $X=S^0$.
\end{definition}

% SOURCE: KIP126/Main/Axiom/Literature/MainPaper/main.tex:126-150 and 253-275; standard mod-2 Adams resolution used implicitly.
\begin{definition}[$H\mathbb F_2$-Adams resolution and tower]
  \label{def:hf2-adams-tower}
  \uses{def:admissible-spectrum,def:hf2-homology,def:cofiber-triangle,def:stable-sphere-suspension}
  \notready
  For an admissible spectrum $X$, iterate the unit
  $X\to H\mathbb F_2\wedge X$ and its fiber to obtain a functorial Adams
  resolution and a decreasing tower
  \[
    \cdots\longrightarrow X_2\longrightarrow X_1\longrightarrow X_0=X.
  \]
  The interface records the exact triangles, augmentations to $X$, naturality,
  external tower pairings
  $X_s\wedge Y_{s'}\to (X\wedge Y)_{s+s'}$, and the fact that every successive
  cofiber is an $H\mathbb F_2$-module spectrum.  The pairings do not give an
  internal multiplication on the tower of an arbitrary $X$.
\end{definition}

% SOURCE: KIP126/Main/Axiom/Literature/MainPaper/main.tex:253-343, Adams filtration on classes and maps.
\begin{definition}[Filtration induced by the Adams tower]
  \label{def:adams-tower-filtration}
  \uses{def:hf2-adams-tower,def:stable-homotopy-group,def:core-filtration}
  \notready
  The Adams filtration on $\pi_nX$ is
  \[
    F^s\pi_nX=\operatorname{im}\bigl(\pi_nX_s\longrightarrow\pi_nX\bigr).
  \]
  A map has filtration at least $k$ when it factors through the $k$th stage of
  the target tower, with the chosen factorization retained as data.  Tower
  pairings give filtration inequalities for external smash products and for
  composition.  They give internal-product inequalities only after a
  multiplication on the spectrum has been supplied.
\end{definition}

% SOURCE: KIP126/Main/Axiom/Literature/MainPaper/main.tex:929-971, especially the homology consequence in
% Proposition prop:41561db2; standard first-stage Adams-tower argument.
\begin{proposition}[Positive Adams filtration is homologically trivial]
  \label{prop:positive-adams-filtration-homology-zero}
  \uses{def:hf2-adams-tower,def:adams-tower-filtration,def:adams-filtration,def:hf2-homology}
  \notready
  For an admissible stable map $u:X\to Y$, if
  $\operatorname{AF}(u)>0$, including infinite filtration, then
  \[
    H_*u=0.
  \]
  Consequently, if $H_*u$ is nonzero then $\operatorname{AF}(u)=0$.  The
  converse direction is not used implicitly: whenever homological vanishing
  is meant to produce positive filtration, a chosen lift through the first
  Adams stage must be supplied separately.
\end{proposition}

\begin{proof}
  Positive filtration supplies a factorization through the first fiber in the
  $H\mathbb F_2$-Adams tower of $Y$.  The composite of that fiber with the
  unit $Y\to H\mathbb F_2\wedge Y$ is null, so applying
  $H\mathbb F_2$-homology makes the induced map zero.  The final assertion is
  its contrapositive, using the nonnegative filtration convention.
  The factorization-level implication is now proved for the constructed
  tower in Proposition~\ref{prop:h6-tower-homology-zero}, with multiplication
  on $H\mathbb F_2$ and preservation of zero maps stated explicitly.
  Its connection to the numerical $\operatorname{AF}$ interface and the
  chosen fixed model remains pending; this full node is not marked complete.
  \uses{prop:h6-tower-homology-zero}
\end{proof}

% SOURCE: KIP126/Main/Axiom/Literature/MainPaper/main.tex:126-150 and standing convergence assumption at 253-263; classical Adams construction.
\begin{theorem}[Adams tower, Ext, and convergence comparison]
  \label{thm:adams-tower-ext-comparison}
  \uses{def:hf2-adams-tower,def:adams-tower-filtration,def:steenrod-ext,thm:filtered-complex-to-spectral-sequence,def:ss-convergence-detection-witness,thm:external-classical-adams-convergence,def:external-result}
  \notready
  The exact couple or filtered cobar complex of the
  $H\mathbb F_2$-Adams tower yields a spectral sequence whose $E_2$ page is
  naturally
  $\operatorname{Ext}_{\mathcal A}^{s,t}
  (H^*(X;\mathbb F_2),\mathbb F_2)$ and whose induced tower filtration is
  Definition~\ref{def:adams-tower-filtration}.  A located classical
  convergence-and-detection witness is an additional input for the admissible
  spectra used in the paper; this theorem does not prove it.
\end{theorem}

\begin{proof}
  Apply the filtered-complex construction to the tower's cobar resolution,
  identify its normalized cochains with the Steenrod cobar complex, and use
  the module--comodule comparison to obtain the bigraded Ext page.  Attach the
  located witness when a later theorem needs the abutment comparison.  Pairings
  are treated separately in Theorem~\ref{thm:adams-external-pairing}; they are
  not silently converted into an internal product for arbitrary $X$.
\end{proof}

% SOURCE: classical Adams tower pairing; KIP126/Main/Axiom/Literature/MainPaper/main.tex:135-159 uses its
% internal specialization only for the sphere Adams spectral sequence.
\begin{theorem}[Adams external pairings and multiplicative specializations]
  \label{thm:adams-external-pairing}
  \uses{thm:adams-tower-ext-comparison,def:hf2-adams-tower,def:unital-multiplication-data,def:multiplicative-ss}
  \notready
  Whenever $X$, $Y$, and $X\wedge Y$ are in the Adams construction, the tower
  pairings induce natural pagewise pairings
  \[
    E_r^{s,t}(X)\otimes E_r^{s',t'}(Y)
      \longrightarrow E_r^{s+s',t+t'}(X\wedge Y)
  \]
  compatible with page passage and satisfying the Leibniz rule.  Taking
  $X=S^0$ makes the sphere Adams spectral sequence an internally unital
  multiplicative spectral sequence, and every $E_r(Y)$ is naturally a module
  over it.  More generally, if $R$ has chosen unital multiplication data, then
  composing the external pairing with $R\wedge R\to R$ gives an internal
  multiplicative Adams spectral sequence for $R$.  No internal multiplication
  is asserted for an arbitrary admissible spectrum.
\end{theorem}

\begin{proof}
  Pass the external pairings of Adams towers to their exact couples, then to
  every page.  Addition of filtration indices gives the displayed bidegree and
  the boundary in the paired filtered complex gives the Leibniz rule.  For the
  sphere, compose with the unit identification
  $S^0\wedge S^0\simeq S^0$; for $R$, compose with its chosen multiplication.
  Associativity and units follow from the supplied algebra coherences, while
  $S^0\wedge Y\simeq Y$ gives the module action.
\end{proof}

% SOURCE: KIP126/Main/Axiom/Literature/MainPaper/main.tex:134-150 and 253-263.
\begin{definition}[Classical Adams spectral sequence]
  \label{def:classical-adams-ss}
  \uses{thm:adams-tower-ext-comparison,thm:adams-external-pairing,def:steenrod-ext,def:adams-bidegree,def:adams-complex-shape,def:ss-elementwise-page-presentation,def:ss-einfty-presentation,def:admissible-spectrum,def:stable-homotopy-group,def:ss-convergence-detection-witness}
  \notready
  For every admissible spectrum $X$, the classical mod--$2$ Adams spectral
  sequence has
  \[
    E_2^{s,t}(X)=\operatorname{Ext}_{\mathcal A}^{s,t}
      (H^*(X;\mathbb F_2),\mathbb F_2),
    \qquad d_r:E_r^{s,t}\to E_r^{s+r,t+r-1},
  \]
  together with the AIM finite-page presentation
  $E_r\cong Z_{r-1}/B_{r-1}$ for $r\ge2$, the algebraic $E_\infty$
  presentation, and naturality.  A chosen convergence-and-detection witness
  identifies this $E_\infty$ with the $2$-complete stable homotopy of $X$ when
  supplied.  These sequences
  have external smash pairings and are modules over the sphere Adams spectral
  sequence.  The sphere sequence, and more generally the sequence of an
  admissible spectrum with chosen unital multiplication data, is internally
  multiplicative; no internal product is included for general $X$.
\end{definition}

% SOURCE: KIP126/Main/Axiom/Literature/MainPaper/main.tex:253-263 and 841-864.
\begin{proposition}[Classical Adams convergence and detection]
  \label{prop:classical-adams-convergence}
  \uses{def:classical-adams-ss,def:ss-convergence-detection-witness,def:adams-detection}
  \notready
  Given a convergence-and-detection witness for the classical sequence, its
  $E_\infty$ page is identified with the associated graded of the Adams
  filtration on the $2$-complete stable homotopy groups.  Consequently the
  witness's detection relation identifies the filtered homotopy classes
  represented by nonzero permanent classes.
\end{proposition}

\begin{proof}
  Project the specified witness to the classical Adams sequence and identify
  its filtration with the Adams-tower filtration.  The algebraic
  $Z_\infty/B_\infty$ presentation then supplies the page representatives used
  by the witness's detection relation.
\end{proof}

% SOURCE: KIP126/Main/Axiom/Literature/MainPaper/main.tex:136-162.
\begin{definition}[The classes $h_j$]
  \label{def:adams-hj}
  \uses{def:classical-adams-ss,thm:adams-external-pairing}
  \notready
  For $j\ge0$, $h_j$ denotes the standard generator in
  $E_2^{1,2^j}(S^0)$.  Its square $h_j^2$ lies in
  $E_2^{2,2^{j+1}}(S^0)$ and has stem $2^{j+1}-2$.  In particular,
  $h_6^2$ lies in bidegree $(2,128)$ and stem $126$.  All products here use
  the canonical internal multiplication on the sphere Adams spectral sequence
  supplied by Theorem~\ref{thm:adams-external-pairing}.
\end{definition}

\section{Adams filtration and detection}

% SOURCE: KIP126/Main/Axiom/Literature/MainPaper/main.tex:253-263, 283-306, and 912-997.
\begin{definition}[Adams filtration of elements and maps]
  \label{def:adams-filtration}
  \uses{def:admissible-spectrum,def:hf2-homology,def:stable-homotopy-group,def:adams-tower-filtration}
  \notready
  The primitive notions are predicates $\operatorname{AF}(\alpha)\ge s$ and
  $\operatorname{AF}(f)\ge k$: respectively, $\alpha$ lies in filtration
  $F^s\pi_nX$, and $f$ factors through the $k$th Adams-tower stage.  Their
  suprema take values in $\mathbb Z\cup\{+\infty\}$ for elements and
  $\mathbb N\cup\{+\infty\}$ for maps; in particular the zero element and zero
  map have infinite filtration.  The paper's notation
  $\operatorname{AF}(-)=k$ abbreviates that the finite supremum is attained
  and equal to $k$.  The interface records monotonicity under sums,
  composition, and the shift induced on detected classes.
\end{definition}

% SOURCE: KIP126/Main/Axiom/Literature/MainPaper/main.tex:307-343, 408-459, and 461-648.
\begin{definition}[Detection relation]
  \label{def:adams-detection}
  \uses{def:ss-convergence-detection-witness,def:adams-filtration}
  \notready
  For a permanent Adams class $x\in E_\infty^{s,t}(X)$ and a homotopy class
  $\alpha\in\pi_{t-s}X$, write $x\leadsto\alpha$ when $x$ is the associated
  graded image of $\alpha$ in filtration $s$.  Detection is a relation rather
  than a function: a homotopy element can have several representatives after
  adding higher-filtration terms.
\end{definition}
